\clearpage\thispagestyle{empty}
\pdfbookmark[0]{Quellen zu Gauß an Schumacher (18. August 1844)}{cometaug18-sources}
\begin{center}\textbf{Quellen zu Gauß an Schumacher (18. August 1844)}\end{center}
\noindent Diese KI-Übertragung ins heutige Deutsch folgt dem Brief auf S.~467--468 der \href{https://gdz.sub.uni-goettingen.de/id/PPN236007467}{\textit{Werke}, Band~VI} (Göttingen, 1874). Carl Friedrich Gauß schrieb am 18. August 1844 an Heinrich Christian Schumacher. Der Erstdruck erschien in der Beilage zu den \textit{Astronomischen Nachrichten}, Band~XXII, Nr.~514, \href{https://api.digitale-sammlungen.de/iiif/image/v2/bsb10049447_00091/full/1800,/0/default.jpg}{Sp.~165--166}. Die \href{https://api.digitale-sammlungen.de/iiif/image/v2/bsb10049447_00092/full/1800,/0/default.jpg}{Schlusszeile auf Sp.~168} datiert die Beilage auf den 5. September 1844. Die Quellenangabe der Werkausgabe nennt das Briefdatum.
\par\smallskip\noindent Der Erstdruck bezeichnet Gauß als Direktor der Göttinger Sternwarte und enthält seine Unterschrift. Die Werkausgabe nennt Schumacher als Empfänger und endet mit dem Hinweis auf den einzigen Vergleich. Ihre Tabelle verwendet hochgestellte Buchstaben für die Zeiteinheiten, Dezimalpunkte und je drei Striche für die fehlenden Winkelangaben. KI-Übertragung und KI-Satz.
\clearpage
